\chapter{Reliability Qualification: Building the Lifetime Confidence Argument}
\label{ch:reliability}

\textit{Read first:} FIT versus DPPM; qualification argument; sample confidence;
burn-in versus HTOL; wear-out families.

\textit{Deep dive:} Arrhenius life projection; GR-468 stress detail; connector
mate-cycle physics.

\textit{Reference:} qualification planning matrix; wear-out map.

\begin{quote}
\noindent\textbf{Reliability qualification at a glance}\\[0.4em]
\begin{itemize}
\item Start with the requirement. Define the life, environment, handling, and
      performance claim that needs support.
\item Identify credible failure mechanisms. Ask what physical or electrical
      process could violate that claim.
\item Select an acceleration method. Choose a stress that meaningfully
      accelerates the mechanism without creating an unrelated failure.
\item Define observables. Decide what measurable change reveals degradation
      before running the stress.
\item Define acceptance criteria. State how much change is allowable and why.
\item Choose representative samples. Cover relevant lots, suppliers, sites,
      process corners, and product variation.
\item Interpret the evidence. Connect observed degradation, confidence, and
      model assumptions to a bounded decision.
\item Feed the result into production controls. Determine what must be
      controlled, screened, sampled, or monitored after qualification
      (\Cref{ch:manufacturing}).
\end{itemize}
Qualification is not evidence because a standard test was performed. It is
evidence when a named stress, observable, sample plan, and acceptance criterion
support a specific reliability decision.
\end{quote}

A five-year lifetime requirement cannot be verified by operating a product for
five years before release. Reliability engineering therefore works indirectly.
The team identifies the physical mechanisms that could cause the product to
degrade, selects stresses that accelerate those mechanisms, measures signatures
that reveal change, and determines whether the resulting evidence supports the
intended use and lifetime.

The stress itself is not the argument. HTOL, humidity, temperature cycling, ESD,
vibration, and connector cycling answer different questions. Running all of them
does not automatically create confidence, and omitting one is not automatically
irresponsible. The qualification plan must connect each activity to a credible
threat.

Qualification also differs from characterization. A temperature sweep during
characterization asks how a healthy product behaves while it is hot or cold. A
qualification stress asks whether exposure causes permanent change. Reversible
movement and lasting degradation require different evidence and different
decisions.

Finally, qualification does not prove that every future production unit will be
reliable. It establishes bounded confidence for the tested design, sample
population, mechanisms, stresses, durations, observables, and assumptions.
\Cref{ch:manufacturing} addresses whether production can reproduce that
qualified population consistently.

\section{What qualification is, and is not}
\label{sec:qual-terms}

Keep these jobs distinct:

\begin{description}
\item[Characterization] Maps behavior and distributions across corners.
\item[Verification] Checks measured results against frozen requirements.
\item[Validation] Demonstrates that the complete product is suitable for its
      intended system use (\Cref{ch:validation}).
\item[Reliability qualification] Builds confidence that named lifetime and
      environmental mechanisms do not violate the intended claim.
\item[Production screening] Attempts to detect unacceptable units or early-life
      defects at manufacturing scale (\Cref{ch:manufacturing}).
\item[Burn-in] A possible production screen for specific early-life mechanisms.
      Burn-in is not automatically qualification and does not replace
      qualification HTOL.
\end{description}

\noindent\textit{Module example.} Characterization maps LIV and BER versus
temperature. Verification checks that ship power and BER meet the PRD.
Validation shows the module works in the target host and plant. Qualification
HTOL asks whether biased high-temperature exposure permanently shifts threshold
or slope beyond the life claim. Burn-in may cull infant mortality on the line;
it does not prove five-year wear-out life.

\section{The language of scale: FIT and DPPM}
\label{sec:fit-dppm}

Fleet arguments need two different numbers. One is about life in the field; the
other is about quality at the factory door. A rate without population
definition, observation time, and confidence is not a fleet claim
(\Cref{sec:tree-evidence-block}).

\begin{description}
\item[FIT (failures in time)]\firstterm{FIT} failures per $10^{9}$ device-hours.
      State the population unit (laser die, module, or link), observed versus
      predicted FIT, confidence interval, useful-life / constant-hazard
      assumption, and whether the system is treated as repairable. Multiply a
      per-unit FIT by population and hours only after those definitions are
      fixed.
\item[DPPM (defective parts per million)] Distinguish incoming defects, outgoing
      defects, and escaped field defects. Always state the denominator, time
      window, and whether rework and retest are included.
\end{description}

\keyidea{Zero failures do not prove a zero failure rate.
0 failures in 20 units is not the same evidence as 0 failures in 1{,}000 units.
0 failures after $10^{3}$ device-hours is not the same evidence as 0 failures
after $10^{6}$ device-hours. Evidence strength depends on sample-hours,
one-sided upper confidence bounds, censoring, and whether lots and sites are
representative. Qualification and production SPC complement each other; neither
alone is a fleet claim.}

\aside{An order-of-magnitude exercise: a
$100{,}000$-GPU cluster with several optical links per GPU has on the order of
$10^{5}$--$10^{6}$ lasers. Even a small per-laser FIT then implies a nontrivial
number of link failures per day, which is exactly why field-replaceable external
laser sources and redundancy are attractive.}

\section{Qualification flows}

Qualification is a life and variation argument, not a list of stresses. A
requirement such as five-year operation is too abstract to test directly. You
identify the mechanisms that could violate it, choose stresses that accelerate
those mechanisms, monitor an observable signature, and decide whether the
resulting evidence is sufficient for the claim
(\Cref{sec:tree-qual-evidence,sec:interview-validation-ladder}).

Debugging is what you do when remaining margin hits zero. Qualification asks
how much margin remains after the expected stresses. Place this work in
\Cref{sec:validation-workflow} Step~6; manufacturing validation and ATP live in
\Cref{ch:manufacturing}. Keep the customer view and the vendor view distinct: the
vendor designs internals; the customer characterizes externally visible behavior
and decides deployment
(\Cref{sec:interview-margin-budget,sec:interview-customer-vendor}).

\begin{dectree}
Requirement
  |
Budget (life / FIT / DPPM / power)
  |
Allocation (die / package / module / host)
  |
Verification (GR-468 / GR-1221 / JESD47 / ATP)
  |
Production + field monitoring
\end{dectree}

\engcheck{Margin budgeting}{%
Every stress spends margin: temperature, voltage, ripple, contamination,
insertion loss, connector wear, vibration, aging, process variation.
Qual verifies what remains, not only that the part still links.}

\engcheck{Customer view vs vendor view}{%
Vendor: internals, device physics, implementation.\\
Customer: BER, sensitivity, FEC, telemetry, environmental sweeps, interop.\\
Engineering samples (Tx-only, Rx-only, breakout, PRBS) open isolation;
otherwise stay on the external surface.}

For a bookended product, begin with BER, FEC, sensitivity, telemetry,
environment, and interop. Request engineering access only when that surface
cannot decide (\Cref{sec:tree-blackbox-access}).

\paragraph{Worked story: gradual laser degradation.}
Walk one mechanism end to end before opening the matrix. The requirement is a
named life claim, for example five years at the use condition and thermal class
in the PRD. The threat is gradual active-region wear: threshold rises and slope
falls over life. That is not COD (sudden dark after a healthy ship LIV) and not
ESD or latch-up (hard fail with supply signature and no optical aging trend).

The evidence strategy is representative lots and sites under justified HTOL:
elevated temperature and bias chosen because they accelerate that mechanism,
with sample size, stress hours, and confidence stated
(\Cref{sec:qual-sample-strategy,sec:fit-dppm}). HTOL is useful only when that
acceleration argument holds. Observables are LIV and system proxies versus the
ship baseline: $I_\mathrm{th}$, slope efficiency, wavelength, launch power or
BER drift, recorded before and during stress at named planes.

Acceptance is a bounded change after projection that still supports the life
claim. Document $E_a$ and confidence when converting HTOL hours to field years,
and do not apply one $E_a$ to mixed mechanisms. Production control is decided last and owned in \Cref{ch:manufacturing}: a
room-temperature LIV/power ATP proxy, a sampled hot audit, SPC on
$I_\mathrm{th}$ and slope, or an explicit statement that no cost-effective 100\%
screen separates the weak tail. A production burn-in may cull infant mortality;
it does not replace this life argument.

Other mechanisms (solder fatigue, contamination, connector wear) follow the same
template. \Cref{tab:qual-planning-matrix} is the reference fill-in, not the
primary explanation.

\paragraph{Qualification planning matrix.}
\label{sec:qual-planning-matrix}

The matrix restates the laser story (and siblings) as
mechanism $\rightarrow$ stress $\rightarrow$ observable $\rightarrow$
acceptance $\rightarrow$ production control. Fill cells for the product class
and claimed life; do not treat blank rows as covered.

\begin{table*}[ht]
\refstepcounter{table}\label{tab:qual-planning-matrix}%
\centering
\setlength{\tabcolsep}{3pt}%
\footnotesize
\renewcommand{\arraystretch}{1.15}%
\begin{tabular}{@{}p{0.18\linewidth}p{0.18\linewidth}p{0.18\linewidth}p{0.18\linewidth}p{0.20\linewidth}@{}}
\toprule
Failure mechanism & Stress & Observable & Acceptance & Production control \\
\midrule
Laser degradation & Temperature / lifetime (HTOL) & $I_\mathrm{th}$, slope, $\lambda$, BER & Named limit vs life claim & ATP / SPC / sampled audit / none \\
\addlinespace[0.7em]
Solder fatigue & Temperature cycling & Resistance, BER, opens & Post-stress continuity / BER & Process control, FAIR \\
\addlinespace[0.7em]
Contamination / corrosion & Humidity / damp heat & Loss, ORL, leakage & IL/ORL / functional limits & Handling, sealing, audit \\
\addlinespace[0.7em]
Connector wear & Mate cycling & Insertion loss, ORL & Cycle-count IL budget & Supplier / hygiene control \\
\bottomrule
\end{tabular}
\par\medskip
{\raggedright\footnotesize\textbf{Table~\thetable.} Qualification planning
matrix (reference). Same template as the laser-degradation story above.
Interview form: \Cref{fw:qual-plan,sec:tree-qual-evidence}. Decision unlocked:
which mechanism-stress-observable row is missing before you claim life.\par}
\end{table*}

\paragraph{GR-468 in practice.}
\label{sec:gr468}

Optoelectronics inherited a common qualification language from telecom:
\term{Telcordia GR-468-CORE}\firstcite{gr468}. Map each stress onto the
qualification evidence path in \Cref{sec:tree-qual-evidence}; do not invent a
second sequence. Core stresses that still show up on every laser and module
program: HTOL for life or mechanism evidence; burn-in as a production
infant-mortality screen when justified; temperature cycling and damp heat; ESD
and mechanical stress.

Keep the jobs distinct: \textbf{burn-in} screens infant mortality from a
production population; \textbf{qualification HTOL} gathers life or mechanism
evidence under accelerated operation. Do not imply that every GR-468-style HTOL
is a per-unit screen. A 1{,}000-hour life test may justify a 100\%
room-temperature proxy, a sampled hot audit, a process monitor, or no direct
production screen at all. Document $E_a$ and confidence bounds when converting
HTOL hours to field years, keep sample-size humility
(\Cref{sec:fit-dppm,ch:manufacturing}), and qualify the laser die, hermetic package, and
module assembly separately when failures split across those boundaries
(\Cref{sec:photonic-packaging,sec:laser-aging,sec:elsfp}).

\term{Arrhenius} acceleration underpins life projection only when the named
failure mechanism is temperature-accelerated in the assumed regime.

\paragraph{Sample strategy and confidence.}
\label{sec:qual-sample-strategy}

Weak answer: ``We test 20 units.'' Strong answer: the sample strategy depends
on the failure-rate target, the confidence requirement, cost, and population
variation. State lots, date codes, suppliers, manufacturing sites or lines,
process corners, and whether the claim is zero-failure upper-bound or observed
rate (\Cref{sec:fit-dppm,sec:tree-evidence-block}). A rate without population,
observation time, and confidence is not a fleet claim.

\heuristic{A small sample that saw zero failures does not prove field life. It
sets an upper bound. State that bound, or do not make the claim.}

\tradeoff{More qualification vs faster release}{Confidence and fewer late
escapes}{Schedule, cost, and delayed learning in the field}{When remaining risk
is high and reversible controls cannot cover it}{Prioritize tests by risk
$\times$ uncertainty $\times$ impact. Do not test everything equally.}

\paragraph{GR-1221: the passive-component companion.}
GR-468 covers active optoelectronics. Its companion, \term{Telcordia GR-1221-CORE}
(Generic Reliability Assurance Requirements for Passive Optical Components), covers
the parts GR-468 does not: connectors, fiber couplers, WDM filters and MUX/DEMUX,
splitters, and isolators~\firstcite{gr1221}. It uses the same style of stress
sequence (damp heat, temperature cycling, mechanical, and aging tests) but scores
pass/fail on insertion loss and return loss rather than on LIV. A short-reach link
that leans on an on-package or blind-mate MUX and on external multi-wavelength
sources carries a passive reliability budget that lives in GR-1221, not GR-468
(\Cref{sec:connector-reliability,ch:wdm}). Split the qual the same way you split
the FIT: active laser die under GR-468, silicon under JESD47, passive optics under
GR-1221.

\paragraph{Handoff to manufacturing.}
Qualification bounds life risk for representative hardware. Whether production
can reproduce that result, screen escapes, and hold SPC is developed in
\Cref{ch:manufacturing}. Failures that pass qual but fail field usually sit in
derating policy, connector contamination, or a manufacturing coverage gap
(\Cref{sec:laser-aging,sec:fleet-triage,sec:escaped-defects}).

\section{Electronics reliability: driver, TIA, and DSP silicon}
\label{sec:ic-reliability}

GR-468 covers the optoelectronic parts of the link: the laser die, the
photodiode, and the hermetic or non-hermetic package around them. The modulator
driver, TIA, retimer, and DSP (\Cref{sec:drivers,sec:pd-tia}) are ordinary CMOS or
SiGe BiCMOS ICs, and they wear out and fail by a different, better-documented set
of mechanisms. Treat them with the semiconductor industry's own qualification
language, not with Arrhenius laser-aging math borrowed from \Cref{sec:laser-aging}.

\paragraph{JESD47: the silicon-side GR-468.}
\firstterm{JESD47} JEDEC JESD47 is the baseline stress-test-driven qualification
flow for a new IC, a device family, or a process change: temperature cycling,
HTOL, HTSL (high-temperature storage life), autoclave or HAST (highly accelerated
stress test) for moisture, and mechanical shock and vibration~\firstcite{jesd47}.
It plays the same role for driver and TIA silicon that GR-468 plays for the
laser: a common list of stresses that a supplier runs once and a customer accepts
instead of renegotiating a qual plan on every design win.

\paragraph{ESD and latch-up: failure modes GR-468 does not test.}
Two mechanisms are specific to ICs and absent from the laser-side wear-out map in
\Cref{tab:wearout-map}:

\begin{description}
\item[ESD] a discharge event during handling or assembly damages a gate oxide or
      junction. Component-level classification uses the human-body model (HBM) and
      charged-device model (CDM) test standards, \term{ANSI/ESDA/JEDEC JS-001} and
      \term{JS-002}~\firstcite{jsesd}. A driver or TIA datasheet HBM/CDM rating is
      the number that protects the part on the factory floor, at fiber-attach and
      wire-bond stations where a laser die is also exposed.
\item[Latch-up] a parasitic thyristor structure in CMOS turns on under an
      overvoltage or current-injection event and holds a low-impedance path until
      power is cycled. \term{JESD78}\firstterm{JESD78} defines the overvoltage and
      $\pm100$~mA current-injection test that classifies susceptibility by supply
      and signal pin~\firstcite{jesd78}. A latched driver IC can look like a dead
      laser on the bench (no light, no LIV signature) until you check the supply
      current instead of the optical path.
\end{description}

Both mechanisms are 100\%-screen or design-margin items, not something you
project with an activation energy. If a driver fails ESD or latch-up in the
field, that is a manufacturability or design-margin bucket item
(\Cref{sec:fleet-triage}), not a wear-out FIT argument.

\paragraph{AEC-Q100: a borrowed grade, not a requirement.}
\term{AEC-Q100}\firstterm{AECQ100} is the automotive industry's qualification
standard for ICs, built on the same JEDEC JESD47/JESD22 stress methods with
tighter ESD targets and named temperature grades from Grade~3
($-40$ to $85$\textdegree C) up to Grade~0 ($-40$ to $150$\textdegree C)~\firstcite{aecq100}.
Datacenter optics does not require Q100; the fleet lives in a controlled data hall,
not an engine bay. It is still a useful signal: a driver, TIA, or retimer die that
also ships in an automotive part number typically carries a published Q100 grade,
and that grade is a fast proxy for the ESD/latch-up margin and temperature-cycle
depth behind the datasheet, without re-running the qual plan yourself.

\paragraph{Where this lands after qualification.}
Carry IC-level qual into the production acceptance and SPC structure in
\Cref{ch:manufacturing}: require the supplier's JESD47 qual report and
HBM/CDM/latch-up ratings for driver and TIA die at DVT, add an ESD handling
audit to the incoming-QC checklist alongside laser LIV/SMSR sampling, and treat
a driver/TIA silicon revision the same way you treat a laser die revision or a
CMIS firmware rev: an ECO that needs first-article requalification, not a silent
BOM swap.

\section{Wear-out modes to know}
\label{sec:wearout-modes}

Arrhenius math, derating, and the worked FIT example live in
\Cref{sec:laser-aging}. This section is the mechanism catalog: how each failure
shows up in ATP and telemetry, and which triage bucket owns it
(\Cref{sec:fleet-triage}). Do not run process CAPA on a wear-out part, and do not
burn FIT math on a dirty connector.

\paragraph{Infant mortality versus wear-out versus packaging.}
Field failures come in three clocks, and mixing them up wastes CAPA. Infant
mortality is early fails from latent defects; burn-in and HTOL screens remove them
before ship (\Cref{sec:gr468}). Wear-out is gradual or sudden end-of-life in the
laser or EAM under temperature, current, and optical power, projected with
Arrhenius and derating (\Cref{sec:laser-aging}). Packaging and assembly faults
(FAU align, epoxy creep, solder voids, connector wear) often dominate field
returns once lasers are screened (\Cref{sec:photonic-packaging}).
Destructive physical analysis (facet cross-section, EDX, FAU section) is required
when the signature is ambiguous or when you need evidence for supplier 8D
(\Cref{sec:supplier-exec}).

\paragraph{Mechanism map.}
\Cref{tab:wearout-map} is the working list for laser-bearing modules and CPO/ELS
paths. Customize limits in the ATP; keep the classification discipline.

\begin{table*}[ht]
\refstepcounter{table}\label{tab:wearout-map}%
\centering
\setlength{\tabcolsep}{3pt}%
\footnotesize
\renewcommand{\arraystretch}{1.25}%
\begin{tabular}{@{}p{0.16\linewidth}p{0.22\linewidth}p{0.28\linewidth}p{0.26\linewidth}@{}}
\toprule
Mechanism & Observable & ATP / telemetry & Triage bucket \\
\midrule
COD (facet) & Sudden dark or hard fail; was healthy & Dark LIV; DPA facet; date-code cluster? & Reliability (COD) or mfg (ESD) \\
\addlinespace[1.5em]
Gradual facet / active region & $I_\mathrm{th}$ up, slope down over life & LIV trend vs ship ATP; HTOL lot history & Reliability (wear-out) \\
\addlinespace[1.5em]
SMSR collapse & Side modes rise; modal noise / BER & OSA SMSR vs floor at $T$ & Reliability; watch aging \\
\addlinespace[1.5em]
EAM aging (EML) & TDECQ/RLM creep at fixed bias & EAM bias sweep + DCA; bias creep log & Reliability (EAM) \\
\addlinespace[1.5em]
RIN rise & BER floor up; feedback sensitive & RIN @ ORL; isolator / connector check & Perf if ORL; reliability if isolator \\
\addlinespace[1.5em]
TEC / thermal control & Unlock or $\lambda$ walk; LIV may look fine & TEC current, case $T$, lock status & Perf (lock) or reliability (TEC) \\
\addlinespace[1.5em]
Coupling / FAU / solder & Loss step, intermittent LOS, shock-related & ORL, mate cycles, DPA FAU/solder & Manufacturability / packaging \\
\addlinespace[1.5em]
Driver/TIA latch-up (ESD) & Sudden hard fail; no light, no LIV signature; supply current spikes & Supply current vs bias; JESD78 rating; date-code cluster? & Mfg (ESD) or design margin \\
\addlinespace[1.5em]
Connector wear / contamination & ORL creep after repeated mate cycles; RIN floor rise & Mate-cycle count vs IEC~61300-2-2 rating; endface grade (IEC~61300-3-35) & Manufacturability / packaging \\
\bottomrule
\end{tabular}
\par\medskip
{\raggedright\footnotesize\textbf{Table~\thetable.} Wear-out and packaging mechanisms
versus observables. Arrhenius projection and derating for the laser rows:
\Cref{sec:laser-aging}. Electronics stress qualification:
\Cref{sec:ic-reliability}. Connector reliability: \Cref{sec:connector-reliability}.
Field classification workflow: \Cref{sec:fleet-triage}.\par}
\end{table*}
\FloatBarrier

\subsection{Reading the wear-out map}
\label{sec:wearout-map-read}

\Cref{tab:wearout-map} is a triage map, not a life calculator. Classify before
corrective action. Reliability rows need life models and derate. Performance
rows need plant and control fixes. Manufacturing and packaging rows need process
and ATP screens. Mixing buckets wastes CAPA (\Cref{sec:fleet-triage}). Arrhenius
projection lives in \Cref{sec:laser-aging}; this section teaches how to read the
map in three families.

\paragraph{Semiconductor wear.}
COD, gradual active-region degradation, SMSR collapse, and EAM aging change the
diode or modulator physics. COD appears as sudden dark after a healthy ship LIV;
gradual wear appears as rising $I_\mathrm{th}$ and falling slope; SMSR collapse
raises modal noise while average power can still look fine; EAM aging creeps
TDECQ/RLM at fixed laser bias. Separate them with dark LIV and facet DPA (COD
versus ESD), LIV trends versus ship ATP (wear), OSA SMSR versus temperature and
age (modal), and EAM bias sweeps with eye metrics (modulator). Controls are life
review and derate for wear, SMSR/aging ATP for modal risk, and EAM life/ATP
screens for modulator aging. Do not burn Arrhenius math on an ESD lot, and do
not replace ``good'' DFBs while the EAM curve walks the eye closed.

\paragraph{Control and electronics.}
TEC/lock faults, driver/TIA latch-up or ESD, and ORL-driven RIN rise change the
environment or the electronics path more often than the laser die. Unlock or
$\lambda$ walk with healthy LIV points at cooler or lock; sudden hard fail with
supply-current spikes and no optical aging signature points at ESD/latch-up; a
BER floor that tracks ORL points at reflections before isolator death. Separate
them with TEC current, case $T$, and lock status; supply current and JESD
ratings with date-code clusters; RIN at stated ORL versus clean plant
(\Cref{sec:rin-values,sec:ic-reliability,sec:lock-validation}). Controls are
lock/thermal policy, ESD handling screens, and connector/reflection budget, not
laser FIT CAPA for every soft floor.

\paragraph{Packaging and optical interfaces.}
Coupling, FAU/solder, and connector wear or contamination produce step loss,
intermittent LOS, shock correlation, or ORL creep after mates. Semiconductor
metrics often stay clean. Separate them with ORL, mate-cycle history, endface
grade (IEC~61300), and DPA of FAU/solder when needed
(\Cref{sec:photonic-packaging,sec:connector-reliability}). Controls are assembly
CAPA, attach screens, service hygiene, and connector ratings. Life models must
not absorb packaging escapes.

Later evidence must not rewrite an earlier wrong bucket without new data. A
clean facet DPA does not clear a connector if ORL was never measured. An HTOL
pass does not clear a date-code ESD cluster.

\section{The reliability bathtub: three failure regimes}
\label{sec:bathtub}

Field failures follow three regimes with different clocks and different fixes.
Mixing them wastes corrective action on the wrong mechanism.

\begin{description}
\item[Infant mortality (early life)] Latent defects from manufacturing: weak
      solder joints, marginal laser die, contamination, firmware bugs that trip on
      first thermal cycle. Burns down rapidly with time. Fixed by burn-in,
      screening, tighter ATP, and first-article control (\Cref{sec:gr468}).
\item[Useful life (constant rate)] Random failures at a roughly steady FIT:
      cosmic-ray-induced single-event upsets, handling damage during unrelated
      service actions, and isolated material defects that passed all screens. Fixed
      by redundancy, field-replaceable modules, and design margin. Note: connector
      wear and contamination accumulate with mate cycles and are usage-driven
      degradation (\Cref{sec:connector-reliability}), not constant-rate random
      failures; budget them separately from the flat-hazard FIT.
\item[Wear-out (end of life)] Physics-driven degradation: laser facet, active
      region, EAM absorption curve shift, TEC aging, epoxy creep. Onset depends
      on temperature, current, optical power, and time. Fixed by derating,
      Arrhenius-based life projection, and planned replacement intervals
      (\Cref{sec:laser-aging}).
\end{description}

A rising failure rate after years of service is wear-out and calls for replacement
planning, not supplier 8D. A cluster of early failures on a new lot is infant
mortality and calls for tighter screens. A steady trickle with no date-code
correlation is useful-life random and calls for redundancy and fast repair. The
triage tree in \Cref{sec:fleet-triage} forces this classification before
corrective action starts.


\section{Photonic packaging and module-level failures}
\label{sec:photonic-packaging}

Fleet FIT is not only laser wear-out. Once lasers are screened and derated,
module and packaging failures often dominate field returns: the part that shipped
with a clean LIV still loses light after shock, humidity, or a thousand ELSFP mate
cycles. Fiber attach and FAU alignment fail from shock, humidity ingress, and
epoxy creep; CPO fiber-array units add assembly steps that wafer test cannot catch
(\Cref{sec:cpo-status}). Hybrid stacks (TFLN-on-Si, InP laser on Si, flip-chip
drivers) introduce solder voids, underfill cracks, and RF return-loss drift
(\Cref{sec:tfln-mzm,sec:drivers}). Thermal paths matter too: uncooled datacom
versus liquid-cooled XPO/CPO, and TEC failure that looks like wavelength drift off
grid or off ring (\Cref{sec:siring,ch:wdm}).

\paragraph{Connector reliability: MPO, mating cycles, and endface quality.}
\label{sec:connector-reliability}

Multi-fiber connectors are the highest-touch mechanical interface in the fleet:
every ELSFP swap, every fiber-attach unit (FAU)\firstterm{FAU} rework, and every cable-plant
install mates and unmates an MPO\firstterm{MPO}. The MPO/MT ferrule family (rectangular, 6.4~mm
$\times$ 2.5~mm, guide-pin aligned, 8/12/16/24 fibers per row) is standardized in
\term{IEC 61754-7}, split into one-fibre-row and two-fibre-row
parts~\firstcite{iec617547}. That standard fixes geometry, not lifetime; lifetime
comes from two companion test methods. \term{IEC 61300-2-2} specifies the
mate/unmate cycling test connector datasheets are rated against, and \term{IEC
61300-3-35} grades endface scratches, pits, and debris into pass/fail zones on the
fiber core and cladding~\firstcite{iec61300}. TIA-568.3 sets 500 cycles as the
structured-cabling
mating-durability floor; MPO/MTP-class connectors in practice are commonly rated
well above 1000 cycles, but that headroom erodes fast with the wrong cleaning
discipline (\Cref{sec:optical-channel}).

Three practical consequences follow for an ELSFP or CPO fiber-attach program.
First, ORL creep is a mating-cycle and cleaning problem before it is a laser
problem: a rising RIN floor after repeated ELS swaps (\Cref{tab:wearout-map}) is
diagnosed with an IEC~61300-3-35-style endface inspection, not a laser FA request.
Second, mate-cycle count belongs in the same telemetry you already read for CMIS
and DDM (\Cref{sec:cmis}); track it per connector, not per module, since a
connector can outlive several module swaps or vice versa. Third, write the
mating-cycle and endface-grade limits into the ATP explicitly
(\Cref{tab:atp-laser}) rather than inheriting a generic MPO datasheet number: an
ELS bank that hot-swaps weekly reaches a 500-cycle floor in under ten years, and a
CPO fiber array that is field-serviced more aggressively reaches it faster still.

ELSFP cycling adds connector wear and contamination that raise ORL
(\Cref{sec:optical-channel,sec:elsfp}); the mating-cycle and endface-grade limits
above are exactly the numbers that turn ``the connector feels loose'' into an ATP
line item instead of a guess.

Destructive physical analysis (cross-section, EDX) and structured 8D/CAPA with
suppliers close the loop from RMA to design rule
(\Cref{sec:supplier-exec,sec:fleet-triage}). Without that loop, packaging FIT gets
mis-attributed to laser Arrhenius models and the wrong part gets redesigned.

\section{From component FIT to fabric availability}
\label{sec:fabric-availability}

The FIT arithmetic in \Cref{sec:fit-example} gives a rate: about $0.6$ laser
failures per day for a fleet of $5\times10^5$ lasers at 50~FIT. That number sizes
the RMA pipeline and the ELS spares bin (\Cref{sec:elsfp}), but it does not say
what a failure costs or how a running job survives one. Two facts turn a
per-component rate into a fabric problem.

First, a training or large inference job is synchronous. A collective
(\Cref{sec:collectives}) waits for its slowest member, so a single dead or slow
link stalls the whole group, not just one endpoint (\Cref{sec:inference-bottlenecks}).
A link that flaps for a second is a stall for every accelerator in that collective.
The optical FIT the earlier chapters budget therefore matters out of proportion to
its share of the parts count.

Second, at cluster scale failures are continuous, not rare. Meta's published
Llama~3 run is the clearest public data point: 16{,}384 H100 GPUs over 54 days
logged 466 interruptions (419 unexpected), roughly one every three hours, while
holding about 90\% effective training time~\firstcite{metallama3}. GPU and HBM3
faults dominated at close to half; network switch and cable faults were 35 events,
8.4\% of the total. The optical link is a minority of hard job stops, but 8.4\% of
a failure every three hours is still tens of network events per run, and the ELS,
module, and connector FIT this chapter budgets
(\Cref{sec:fit-dppm,sec:connector-reliability}) lands in exactly that bucket.

So the design question shifts. It is no longer ``how reliable is one link'' but
``how does a fabric of $10^5$ links keep a job running through a failure every few
hours.'' The answers are architectural, and the optical engineer feeds each one.

\begin{description}
\item[Redundancy and rails.] Rail-optimized topologies (\Cref{sec:topologies})
      already give parallel planes; dual-plane and dual-ToR designs let a lost link
      degrade bandwidth instead of dropping an endpoint. Redundancy multiplies the
      link and laser count, which feeds straight back into the FIT budget: more
      resilience is more parts that can fail.
\item[Detection and reroute.] Transient faults stay below the job. KP4 FEC
      (\Cref{sec:kp4}) absorbs the error bursts a marginal link throws; link-level
      retry and sub-second link-flap detection plus adaptive routing steer traffic
      off a degraded link before the scheduler notices. Vendor fabrics (NVIDIA
      Spectrum-X and Quantum, Broadcom Tomahawk) advertise adaptive or ``cognitive''
      routing and link-level retry for this. Treat the specifics as vendor
      orientation, but the mechanism is why transient optical faults rarely reach
      the hard-stop bucket above.
\item[Topology reconfiguration.] When a link or rack dies for good, an optical
      circuit switch re-wires the topology around it in milliseconds, so the
      scheduler routes around the dead node instead of stalling the pod
      (\Cref{sec:ocs})~\citep{tpuv4ocs}. Component FIT still applies; the fabric
      survives each failure by re-wiring optically.
\item[Sparing and field service.] Hot spare nodes and lanes cover the interval
      between failure and repair. Field-replaceable external lasers (\Cref{sec:elsfp})
      make a dead laser a faceplate swap rather than a fabric outage, which is the
      architectural reason ELS decouples laser FIT from switch FIT. The connector
      mating-cycle and endface budget (\Cref{sec:connector-reliability}) sets how
      many of those swaps the plant survives.
\end{description}

The cost of a failure closes the loop. A hard interruption is lost compute plus
the time to detect, reroute or reschedule, and restart from the last checkpoint.
Fast detection and reroute shrink that lost time, which is the fabric-level reason
the module work in this chapter pays off: derating (\Cref{sec:laser-aging}),
burn-in and screens (\Cref{sec:gr468}), and a tight ATP (\Cref{sec:supplier-exec})
lower the failure rate, and a resilient fabric lowers the cost of each failure that
slips through. The two multiply.
\section{Engineering lens}
\label{sec:reliability-engineering-lens}

\subsection{How it works}

At fleet scale, a modest per-part FIT times millions of parts is a steady stream
of failures. Qualification, derating, and mechanism classification keep that
stream small and classifiable. Manufacturing validation
(\Cref{ch:manufacturing}) then asks whether the factory can reproduce the
qualified population.

\subsection{How it is measured}
\label{sec:reliability-how-measured}

Reliability is measured as a distribution over stress and time. Qualification
records failures and drift by mechanism, stress, lot, and sample history.
Convert accelerated hours to field years only with a named mechanism, stated
$E_a$, and confidence bounds (\Cref{sec:fit-dppm,sec:gr468}). Fleet data add
install age, temperature, firmware, supplier lot, and return code so wear-out
and infant mortality are not mixed (\Cref{sec:fleet-triage}).

\subsection{How it fails}
\label{sec:reliability-how-fails}

Programs fail at life through wrong mechanisms, weak samples, and overstated
acceleration. A stress without a mechanism is only exposure. Zero fails in a
small lot is an upper bound, not a FIT claim. Arrhenius math on an ESD cluster
or dirty connector wastes CAPA.

\subsection{How it is debugged}
\label{sec:reliability-how-debugged}

Classify the bathtub regime and wear-out map bucket before corrective action
(\Cref{sec:bathtub,tab:wearout-map}). Confirm the mechanism with LIV, supply
current, ORL, and DPA as needed. Escalate supplier process CAPA only when the
bucket is manufacturing; escalate life model and derate when the bucket is
wear-out.

\section{Interview takeaway}
\label{sec:reliability-design-review}

\keyidea{Reliability qualification is a mechanism-driven confidence argument.
Begin with the lifetime or environmental claim, identify how it could be
violated, select a credible acceleration method, define observable degradation,
and choose representative samples and acceptance criteria before running the
stress. The result supports a bounded decision; it does not guarantee every
production unit. \Cref{ch:manufacturing} explains how production reproduces and
controls the qualified design.}

Junior mistake: treat zero fails in a small HTOL lot as a FIT claim, or treat a
temperature sweep as temperature qualification
(\Cref{sec:qual-sample-strategy,sec:qual-terms}).

\subsection{Interview Q\&A: Reliability Qualification}
\label{sec:reliability-interview-qa}

Practice speaking these answers aloud. Prefer first-person reasoning over
stress inventories. Planning detail lives in
\Cref{sec:qual-planning-matrix,tab:qual-planning-matrix,sec:qual-terms}.

\paragraph{Question 1. What is reliability qualification?}

\textit{Strong interview answer.}
``Reliability qualification is a bounded confidence argument that a design can
continue meeting its requirements through its intended life and environmental
exposure. I start with the life claim, identify credible failure mechanisms,
select stresses that accelerate those mechanisms, define observable degradation
and acceptance criteria, and choose representative samples. The stress itself is
not the evidence. The evidence is the connection between the requirement,
mechanism, stress, measurement, result, and decision.''

\textit{Likely follow-up.} How is this different from validation? Does passing a
standard mean the product is qualified?

\textit{What the interviewer is testing.} Whether you treat qualification as a
mechanism-driven confidence argument, not a stress checklist.

\textit{Common weak answer.}
``Qualification means running HTOL, humidity, temperature cycling, and
vibration.''

\paragraph{Question 2. What is the difference between validation and
qualification?}

\textit{Strong interview answer.}
``Validation asks whether the product is suitable for its intended system use.
It includes behavior, margin, interoperability, and system corners.
Qualification asks whether time, handling, or environmental exposure creates
permanent degradation that threatens the intended claim. A temperature sweep may
validate operation at high temperature, while repeated cycling or high-temperature
life testing may qualify against fatigue or aging. Similar equipment can be used,
but the questions and acceptance criteria differ''
(\Cref{ch:validation,sec:qual-terms}).

\textit{Likely follow-up.} Is manufacturing validation another form of
qualification? Can characterization data support qualification?

\textit{What the interviewer is testing.} Whether you keep temporary operating
behavior distinct from permanent degradation evidence.

\textit{Common weak answer.}
``Validation is before qualification, and qualification is more severe.''

\paragraph{Question 3. How do you design a qualification plan?}

\textit{Strong interview answer.}
``I start with the lifetime and environmental requirements. Then I build a
failure-mechanism hypothesis space: laser degradation, solder fatigue,
humidity-driven corrosion, ESD damage, connector wear, control-system
degradation, and so on. For each credible mechanism, I choose an acceleration
method, define what measurable signature would reveal degradation, and set
acceptance criteria before running the test. I then choose samples that represent
relevant lots, suppliers, sites, and process corners. Finally, I define how each
result affects release, derating, redesign, production controls, or additional
evidence'' (\Cref{tab:qual-planning-matrix}).

\textit{Likely follow-up.} How do you prioritize mechanisms? What if the standard
requires a test you consider low risk?

\textit{What the interviewer is testing.} Whether you design from mechanisms and
decisions, not from a default stress matrix.

\textit{Common weak answer.}
``I would follow GR-468 and run the required stress matrix.''

\paragraph{Question 4. What is the difference between a failure mode and a
failure mechanism?}

\textit{Strong interview answer.}
``The failure mode is the observable symptom. For example, low optical power,
high BER, or a dead lane. The failure mechanism is the physical or electrical
process that caused it, such as active-region degradation, solder cracking,
corrosion, contamination, or electromigration. Qualification must be designed
around mechanisms, because several mechanisms can produce the same mode and
require different stresses and controls'' (\Cref{tab:wearout-map}).

\textit{Likely follow-up.} Give two mechanisms that could cause low optical
power. Why does this distinction matter to ATP?

\textit{What the interviewer is testing.} Whether you can separate symptom from
physics before choosing stress or CAPA.

\textit{Common weak answer.}
``The failure mode is the cause, and the mechanism is how the unit fails.''

\paragraph{Question 5. Why does accelerated testing predict field life?}

\textit{Strong interview answer.}
``It predicts field life only when the stress accelerates the same failure
mechanism that is expected in use. For temperature-activated degradation, for
example, an Arrhenius relationship may translate elevated-temperature exposure
into an acceleration factor. But the activation energy and model assumptions must
be relevant to the actual mechanism. If the stress creates a different failure,
the test may be severe without being predictive. I would always state the assumed
physics, acceleration model, observables, and uncertainty''
(\Cref{sec:gr468,sec:laser-aging}).

\textit{Likely follow-up.} What is the Arrhenius relationship? How do you know
the activation energy?

\textit{What the interviewer is testing.} Whether you treat acceleration as
mechanism-specific, not universal severity.

\textit{Common weak answer.}
``Higher temperature makes everything age faster, so we can convert hours
directly into years.''

\paragraph{Question 6. Explain HTOL and what you would monitor.}

\textit{Strong interview answer.}
``High-temperature operating life uses elevated temperature while the device is
electrically active to accelerate mechanisms associated with biased operation,
such as laser or semiconductor degradation. I would establish trusted pre-stress
baselines and monitor parameters that reveal gradual movement, not only
catastrophic failure. For a transmitter, that may include threshold current,
slope efficiency, optical power, OMA, wavelength, SMSR, RIN, module current, BER,
and control-loop headroom. I would include intermediate read points because a
unit can remain functional while drifting steadily toward a later cliff.''

\textit{Likely follow-up.} Why is chamber temperature not enough? How would you
choose stress temperature and duration?

\textit{What the interviewer is testing.} Whether you monitor degradation trends
and relevant temperatures, not only end-of-test function.

\textit{Common weak answer.}
``HTOL means running the module hot for a thousand hours and checking whether it
still works.''

\paragraph{Question 7. What does temperature cycling qualify?}

\textit{Strong interview answer.}
``Temperature cycling primarily targets mechanical and package mechanisms caused
by repeated expansion and contraction. Different materials have different
coefficients of thermal expansion, so repeated cycles can fatigue solder joints,
wire bonds, adhesives, fiber attach, and optical alignment. I would monitor
electrical continuity, intermittent lane behavior, optical power, sensitivity,
BER, and permanent alignment movement. The key distinction is whether the change
is reversible with temperature or remains after the cycle.''

\textit{Likely follow-up.} How is thermal shock different? Would you measure
during or only after cycling?

\textit{What the interviewer is testing.} Whether you connect cycling to package
mechanics, not to hot/cold operation alone.

\textit{Common weak answer.}
``Temperature cycling checks that the module operates when it is hot and cold.''

\paragraph{Question 8. What is damp-heat or humidity testing trying to reveal?}

\textit{Strong interview answer.}
``Humidity testing targets moisture-related mechanisms such as corrosion,
leakage, material degradation, contamination movement, and loss of package or
interface integrity. Depending on the mechanism, the unit may be biased or
unbiased during exposure. I would inspect electrical leakage, optical loss,
return loss, power, BER, visual evidence, and selected destructive analysis. A
post-test functional pass alone may miss corrosion or degradation that has begun
but has not yet crossed the system limit.''

\textit{Likely follow-up.} Why might bias matter during humidity testing? How
would you distinguish contamination from corrosion?

\textit{What the interviewer is testing.} Whether you look for moisture
mechanisms beyond waterproof packaging slogans.

\textit{Common weak answer.}
``Humidity testing checks whether the package is waterproof.''

\paragraph{Question 9. Why perform ESD qualification if production units already
pass functional test?}

\textit{Strong interview answer.}
``ESD can cause immediate failure, latent damage, or lost margin without a
complete functional failure. A unit may still link but show higher leakage, worse
receiver sensitivity, increased power, or reduced robustness. Qualification
establishes whether the design and protection network tolerate the intended
handling exposure. Production controls then address grounding, handling,
packaging, and process compliance. A final functional test does not necessarily
detect latent ESD damage'' (\Cref{sec:ic-reliability}).

\textit{Likely follow-up.} What parameters would you compare before and after
ESD? Can ATP screen latent ESD damage reliably?

\textit{What the interviewer is testing.} Whether you understand latent damage and
the split between design qual and handling controls.

\textit{Common weak answer.}
``ESD is required because operators may touch the board.''

\paragraph{Question 10. How do you choose sample size for qualification?}

\textit{Strong interview answer.}
``There is no universal sample count. I choose the plan based on the failure-rate
or degradation claim, required confidence, stress duration, acceleration,
mechanism variability, and population diversity. I also need relevant lots,
suppliers, process corners, and sites. Twenty units with zero failures do not
prove zero field failures. I would report the confidence bound, sample-hours,
censoring, model assumptions, and remaining uncertainty rather than treating a
zero-failure result as absolute proof''
(\Cref{sec:fit-dppm,sec:qual-sample-strategy}).

\textit{Likely follow-up.} What does zero failures actually tell you? Is lot
diversity more valuable than adding several units from the same lot?

\textit{What the interviewer is testing.} Statistical humility and
representativeness, not a folklore sample count.

\textit{Common weak answer.}
``We normally use twenty or thirty units because that is standard.''

\paragraph{Question 11. How do you know when qualification evidence is
sufficient?}

\textit{Strong interview answer.}
``The evidence is sufficient when it supports the intended release decision with
acceptable remaining risk. I look at mechanism coverage, sample
representativeness, degradation distributions, confidence, acceleration-model
uncertainty, and whether failures or trends remain unexplained. The evidence
required for a reversible limited pilot is lower than for unrestricted
deployment. I would not claim that qualification removes uncertainty; I would
state what claim is supported, under which conditions, and what fleet or
production controls remain.''

\textit{Likely follow-up.} Who should accept the remaining risk? What would
justify qualification with restrictions?

\textit{What the interviewer is testing.} Decision-scaled evidence, not
zero-uncertainty fantasies.

\textit{Common weak answer.}
``Qualification is complete when all samples pass all required stresses.''

\paragraph{Question 12. What should you do when a qualification test fails?}

\textit{Strong interview answer.}
``First, I verify the test setup and preserve the unit state. Then I scope
whether the failure is isolated, lot-correlated, stress-dependent, or
systematic. I distinguish the observed failure mode from the suspected mechanism
and use discriminating measurements, controlled swaps, physical analysis, or
reproduction to confirm it. I contain the affected population, determine whether
the stress was relevant, and decide whether the correct response is design
change, process change, supplier control, derating, or additional evidence. After
corrective action, I repeat enough qualification to demonstrate that the
mechanism has been addressed.''

\textit{Likely follow-up.} Would you immediately restart the entire qualification
matrix? What if the failure occurs outside the intended use condition?

\textit{What the interviewer is testing.} Containment, confirmation, and scoped
requal, not panic retesting.

\textit{Common weak answer.}
``I would send the unit to failure analysis and rerun the test.''

\paragraph{Question 13. What is the difference between qualification, burn-in,
and production screening?}

\textit{Strong interview answer.}
``Qualification uses representative samples and accelerated stresses to support
design-life and environmental claims. Burn-in is a possible production screen
intended to precipitate specific early-life defects before shipment. Production
screening includes faster tests or controls applied to every unit or selected
production samples. Burn-in does not replace qualification, and qualification
does not prove every shipped unit was assembled correctly. Each addresses a
different part of the risk'' (\Cref{sec:qual-terms,ch:manufacturing}).

\textit{Likely follow-up.} When would burn-in be justified? What are the
downsides of burn-in?

\textit{What the interviewer is testing.} Whether you keep life evidence distinct
from production screens.

\textit{Common weak answer.}
``Burn-in is qualification performed on every unit.''

\paragraph{Question 14. Give an example of a complete qualification argument for
laser aging.}

\textit{Strong interview answer.}
``Suppose the requirement is five years of continuous operation across the stated
case-temperature range. One credible threat is active-region degradation causing
higher threshold current, lower slope efficiency, reduced OMA, or increased BER.
I would select representative laser and module lots, choose a biased
high-temperature stress using a justified acceleration model, and establish
baseline LIV, wavelength, SMSR, RIN, OMA, module power, and BER. I would define
acceptable drift based on the remaining system margin before the test.
Intermediate read points would show whether degradation is stable or
accelerating. The result would support release, derating, supplier restriction,
thermal redesign, or additional evidence. Production might then monitor
bias-current distributions or sampled LIV, but that production control would not
replace the qualification argument.''

\textit{Likely follow-up.} Why monitor threshold current? How would you
distinguish laser degradation from coupling movement?

\textit{What the interviewer is testing.} End-to-end mechanism argument with
acceptance and downstream control.

\textit{Common weak answer.}
``I would run HTOL and make sure optical power remains in specification.''

\paragraph{Question 15. Give me a 60-second answer for how you would qualify a
new optical transceiver.}

\textit{Strong interview answer.}
``I would begin with the lifetime, environmental, handling, and performance
claims. Then I would identify credible mechanisms that could violate them, for
example laser degradation, solder fatigue, humidity-related corrosion, ESD
damage, connector wear, and package movement. For each mechanism I would select
an acceleration method, define observable degradation and acceptance criteria,
and choose representative samples across lots, suppliers, sites, and process
variation. During stress I would monitor trends rather than only catastrophic
failures. I would interpret the results using appropriate confidence and
acceleration assumptions, investigate any failures to a confirmed mechanism, and
translate the findings into release conditions, derating, design changes,
supplier controls, or manufacturing monitoring. I would communicate the supported
claim and the remaining risk explicitly.''

\textit{Likely follow-up.} Which test would you prioritize if schedule were
limited? What evidence would cause you to hold release?

\textit{What the interviewer is testing.} Staff-level prioritization under
incomplete information.

\textit{Common weak answer.}
``I would follow the qualification standard and run HTOL, humidity, cycling, ESD,
and vibration.''

\subsection{Self-assessment rubric}
\label{sec:reliability-interview-rubric}

Score each spoken answer from 0 to 2 on every dimension. Maximum score is 10.

\begin{description}
\item[Requirement and scope]
      0: no requirement or context.\\
      1: mentions the requirement.\\
      2: defines the requirement, reference plane, condition, and release
      decision.
\item[Mechanism reasoning]
      0: jumps to one cause or lists tests.\\
      1: mentions several possibilities.\\
      2: builds a credible mechanism-based hypothesis space.
\item[Evidence selection]
      0: lists instruments.\\
      1: chooses a useful measurement.\\
      2: explains why the measurement separates hypotheses or supports the
      decision.
\item[Decision quality]
      0: no decision.\\
      1: technical conclusion only.\\
      2: includes release, containment, restriction, redesign, or control.
\item[Communication]
      0: rambling or jargon-heavy.\\
      1: correct but difficult to follow.\\
      2: direct, structured, and appropriately qualified.
\end{description}

Interpretation: 0--4, review the chapter concepts; 5--7, technically competent
but needs stronger structure; 8--9, strong interview response; 10, Staff-level
reasoning and communication. Habit to keep: requirement $\rightarrow$ mechanism
$\rightarrow$ evidence $\rightarrow$ confidence $\rightarrow$ decision
$\rightarrow$ control.
